% Options for packages loaded elsewhere
\PassOptionsToPackage{unicode}{hyperref}
\PassOptionsToPackage{hyphens}{url}
%
\documentclass[
]{article}
\usepackage{amsmath,amssymb}
\usepackage{lmodern}
\usepackage{iftex}
\ifPDFTeX
  \usepackage[T1]{fontenc}
  \usepackage[utf8]{inputenc}
  \usepackage{textcomp} % provide euro and other symbols
\else % if luatex or xetex
  \usepackage{unicode-math}
  \defaultfontfeatures{Scale=MatchLowercase}
  \defaultfontfeatures[\rmfamily]{Ligatures=TeX,Scale=1}
\fi
% Use upquote if available, for straight quotes in verbatim environments
\IfFileExists{upquote.sty}{\usepackage{upquote}}{}
\IfFileExists{microtype.sty}{% use microtype if available
  \usepackage[]{microtype}
  \UseMicrotypeSet[protrusion]{basicmath} % disable protrusion for tt fonts
}{}
\makeatletter
\@ifundefined{KOMAClassName}{% if non-KOMA class
  \IfFileExists{parskip.sty}{%
    \usepackage{parskip}
  }{% else
    \setlength{\parindent}{0pt}
    \setlength{\parskip}{6pt plus 2pt minus 1pt}}
}{% if KOMA class
  \KOMAoptions{parskip=half}}
\makeatother
\usepackage{xcolor}
\usepackage[margin=1in]{geometry}
\usepackage{color}
\usepackage{fancyvrb}
\newcommand{\VerbBar}{|}
\newcommand{\VERB}{\Verb[commandchars=\\\{\}]}
\DefineVerbatimEnvironment{Highlighting}{Verbatim}{commandchars=\\\{\}}
% Add ',fontsize=\small' for more characters per line
\usepackage{framed}
\definecolor{shadecolor}{RGB}{248,248,248}
\newenvironment{Shaded}{\begin{snugshade}}{\end{snugshade}}
\newcommand{\AlertTok}[1]{\textcolor[rgb]{0.94,0.16,0.16}{#1}}
\newcommand{\AnnotationTok}[1]{\textcolor[rgb]{0.56,0.35,0.01}{\textbf{\textit{#1}}}}
\newcommand{\AttributeTok}[1]{\textcolor[rgb]{0.77,0.63,0.00}{#1}}
\newcommand{\BaseNTok}[1]{\textcolor[rgb]{0.00,0.00,0.81}{#1}}
\newcommand{\BuiltInTok}[1]{#1}
\newcommand{\CharTok}[1]{\textcolor[rgb]{0.31,0.60,0.02}{#1}}
\newcommand{\CommentTok}[1]{\textcolor[rgb]{0.56,0.35,0.01}{\textit{#1}}}
\newcommand{\CommentVarTok}[1]{\textcolor[rgb]{0.56,0.35,0.01}{\textbf{\textit{#1}}}}
\newcommand{\ConstantTok}[1]{\textcolor[rgb]{0.00,0.00,0.00}{#1}}
\newcommand{\ControlFlowTok}[1]{\textcolor[rgb]{0.13,0.29,0.53}{\textbf{#1}}}
\newcommand{\DataTypeTok}[1]{\textcolor[rgb]{0.13,0.29,0.53}{#1}}
\newcommand{\DecValTok}[1]{\textcolor[rgb]{0.00,0.00,0.81}{#1}}
\newcommand{\DocumentationTok}[1]{\textcolor[rgb]{0.56,0.35,0.01}{\textbf{\textit{#1}}}}
\newcommand{\ErrorTok}[1]{\textcolor[rgb]{0.64,0.00,0.00}{\textbf{#1}}}
\newcommand{\ExtensionTok}[1]{#1}
\newcommand{\FloatTok}[1]{\textcolor[rgb]{0.00,0.00,0.81}{#1}}
\newcommand{\FunctionTok}[1]{\textcolor[rgb]{0.00,0.00,0.00}{#1}}
\newcommand{\ImportTok}[1]{#1}
\newcommand{\InformationTok}[1]{\textcolor[rgb]{0.56,0.35,0.01}{\textbf{\textit{#1}}}}
\newcommand{\KeywordTok}[1]{\textcolor[rgb]{0.13,0.29,0.53}{\textbf{#1}}}
\newcommand{\NormalTok}[1]{#1}
\newcommand{\OperatorTok}[1]{\textcolor[rgb]{0.81,0.36,0.00}{\textbf{#1}}}
\newcommand{\OtherTok}[1]{\textcolor[rgb]{0.56,0.35,0.01}{#1}}
\newcommand{\PreprocessorTok}[1]{\textcolor[rgb]{0.56,0.35,0.01}{\textit{#1}}}
\newcommand{\RegionMarkerTok}[1]{#1}
\newcommand{\SpecialCharTok}[1]{\textcolor[rgb]{0.00,0.00,0.00}{#1}}
\newcommand{\SpecialStringTok}[1]{\textcolor[rgb]{0.31,0.60,0.02}{#1}}
\newcommand{\StringTok}[1]{\textcolor[rgb]{0.31,0.60,0.02}{#1}}
\newcommand{\VariableTok}[1]{\textcolor[rgb]{0.00,0.00,0.00}{#1}}
\newcommand{\VerbatimStringTok}[1]{\textcolor[rgb]{0.31,0.60,0.02}{#1}}
\newcommand{\WarningTok}[1]{\textcolor[rgb]{0.56,0.35,0.01}{\textbf{\textit{#1}}}}
\usepackage{longtable,booktabs,array}
\usepackage{calc} % for calculating minipage widths
% Correct order of tables after \paragraph or \subparagraph
\usepackage{etoolbox}
\makeatletter
\patchcmd\longtable{\par}{\if@noskipsec\mbox{}\fi\par}{}{}
\makeatother
% Allow footnotes in longtable head/foot
\IfFileExists{footnotehyper.sty}{\usepackage{footnotehyper}}{\usepackage{footnote}}
\makesavenoteenv{longtable}
\usepackage{graphicx}
\makeatletter
\def\maxwidth{\ifdim\Gin@nat@width>\linewidth\linewidth\else\Gin@nat@width\fi}
\def\maxheight{\ifdim\Gin@nat@height>\textheight\textheight\else\Gin@nat@height\fi}
\makeatother
% Scale images if necessary, so that they will not overflow the page
% margins by default, and it is still possible to overwrite the defaults
% using explicit options in \includegraphics[width, height, ...]{}
\setkeys{Gin}{width=\maxwidth,height=\maxheight,keepaspectratio}
% Set default figure placement to htbp
\makeatletter
\def\fps@figure{htbp}
\makeatother
\setlength{\emergencystretch}{3em} % prevent overfull lines
\providecommand{\tightlist}{%
  \setlength{\itemsep}{0pt}\setlength{\parskip}{0pt}}
\setcounter{secnumdepth}{-\maxdimen} % remove section numbering
\ifLuaTeX
  \usepackage{selnolig}  % disable illegal ligatures
\fi
\IfFileExists{bookmark.sty}{\usepackage{bookmark}}{\usepackage{hyperref}}
\IfFileExists{xurl.sty}{\usepackage{xurl}}{} % add URL line breaks if available
\urlstyle{same} % disable monospaced font for URLs
\hypersetup{
  pdftitle={Statistique inférentielle},
  pdfauthor={Vincent Le Meur},
  hidelinks,
  pdfcreator={LaTeX via pandoc}}

\title{Statistique inférentielle}
\author{Vincent Le Meur}
\date{2023}

\begin{document}
\maketitle

\hypertarget{tp2-intervalles-de-confiance}{%
\section{TP2 : Intervalles de
confiance}\label{tp2-intervalles-de-confiance}}

\hypertarget{compte-rendu}{%
\subsection{Compte rendu}\label{compte-rendu}}

\hypertarget{question-1}{%
\subsubsection{Question 1}\label{question-1}}

Calculer la moyenne (empirique) et la variance (empirique) du poids à la
naissance pour l'ensemble des nouveau-nés.

Dans un premier temps on importe le dossier poids.csv

\begin{Shaded}
\begin{Highlighting}[]
\FunctionTok{library}\NormalTok{(readr)}
\NormalTok{poids }\OtherTok{\textless{}{-}} \FunctionTok{read\_csv}\NormalTok{(}\StringTok{"poids.csv"}\NormalTok{, }\AttributeTok{col\_names =} \ConstantTok{FALSE}\NormalTok{)}
\end{Highlighting}
\end{Shaded}

\begin{verbatim}
## Rows: 9863 Columns: 1
## -- Column specification --------------------------------------------------------
## Delimiter: ","
## dbl (1): X1
## 
## i Use `spec()` to retrieve the full column specification for this data.
## i Specify the column types or set `show_col_types = FALSE` to quiet this message.
\end{verbatim}

\begin{Shaded}
\begin{Highlighting}[]
\FunctionTok{View}\NormalTok{(poids)}
\end{Highlighting}
\end{Shaded}

On rappelle les formules de la moyenne empirique et de la variance
empirique:

Moyenne empirique: \[(\frac{1}{n}) \sum(log(X))\]

Variance empirique: \[(\frac{1}{(n-1)}*\sum((log(X)-moy-emp)^2))\]

\begin{Shaded}
\begin{Highlighting}[]
\NormalTok{n}\OtherTok{=}\FunctionTok{length}\NormalTok{(poids}\SpecialCharTok{$}\NormalTok{X1)}


\NormalTok{moy\_emp}\OtherTok{=}\NormalTok{ (}\DecValTok{1}\SpecialCharTok{/}\NormalTok{n)}\SpecialCharTok{*}\FunctionTok{sum}\NormalTok{(poids}\SpecialCharTok{$}\NormalTok{X1)}
\NormalTok{var\_emp}\OtherTok{=}\NormalTok{ (}\DecValTok{1}\SpecialCharTok{/}\NormalTok{n}\DecValTok{{-}1}\NormalTok{)}\SpecialCharTok{*}\FunctionTok{sum}\NormalTok{(((poids}\SpecialCharTok{$}\NormalTok{X1)}\SpecialCharTok{{-}}\FunctionTok{mean}\NormalTok{(poids}\SpecialCharTok{$}\NormalTok{X1))}\SpecialCharTok{\^{}}\DecValTok{2}\NormalTok{)}

\FunctionTok{library}\NormalTok{(knitr)}

\CommentTok{\# Créer le tableau}
\NormalTok{tableau }\OtherTok{\textless{}{-}} \FunctionTok{data.frame}\NormalTok{(}
\NormalTok{  Donné }\OtherTok{=} \FunctionTok{c}\NormalTok{(}\StringTok{"Moyenne empirique"}\NormalTok{, }\StringTok{"Variance empirique"}\NormalTok{),}
  \AttributeTok{Valeur =} \FunctionTok{c}\NormalTok{(moy\_emp, var\_emp)}
\NormalTok{)}

\CommentTok{\# Afficher le tableau avec kable}
\FunctionTok{kable}\NormalTok{(tableau)}
\end{Highlighting}
\end{Shaded}

\begin{longtable}[]{@{}lr@{}}
\toprule()
Donné & Valeur \\
\midrule()
\endhead
Moyenne empirique & 3.125604e+03 \\
Variance empirique & -1.887716e+09 \\
\bottomrule()
\end{longtable}

\hypertarget{question-2}{%
\subsubsection{Question 2}\label{question-2}}

Écrire une fonction qui, à partir d'un vecteur de données et d'un entier
n, renvoie une liste contenant un échantillon aléatoire de taille n
provenant des données fournies, ainsi que la moyenne empirique et la
variance empirique de cet échantillon. Estimer à l'aide de cette
fonction la moyenne et la variance du poids à la naissance à partir des
échantillons de taille n = 30, 60 et 80.

\begin{Shaded}
\begin{Highlighting}[]
\NormalTok{analyse }\OtherTok{\textless{}{-}} \ControlFlowTok{function}\NormalTok{(X,n)\{}
\NormalTok{  echantillon }\OtherTok{\textless{}{-}} \FunctionTok{sample}\NormalTok{(X,n)}
\NormalTok{  moy\_emp}\OtherTok{=}\NormalTok{ (}\DecValTok{1}\SpecialCharTok{/}\NormalTok{n)}\SpecialCharTok{*}\FunctionTok{sum}\NormalTok{(echantillon)}
\NormalTok{  var\_emp}\OtherTok{=}\NormalTok{ (}\DecValTok{1}\SpecialCharTok{/}\NormalTok{n}\DecValTok{{-}1}\NormalTok{)}\SpecialCharTok{*}\FunctionTok{sum}\NormalTok{(((echantillon)}\SpecialCharTok{{-}}\FunctionTok{mean}\NormalTok{(echantillon))}\SpecialCharTok{\^{}}\DecValTok{2}\NormalTok{)}
  
  \FunctionTok{return}\NormalTok{(}\FunctionTok{list}\NormalTok{(}\AttributeTok{echantillon=}\NormalTok{echantillon,}\AttributeTok{moy\_emp=}\NormalTok{moy\_emp,}\AttributeTok{var\_emp=}\NormalTok{var\_emp))}
\NormalTok{\}}
\end{Highlighting}
\end{Shaded}

On Estime maintenant avec différentes tailles d'échantillons:

Pour une taille d'échantillons de 30:

\begin{Shaded}
\begin{Highlighting}[]
\NormalTok{ech30}\OtherTok{=}\FunctionTok{analyse}\NormalTok{(poids}\SpecialCharTok{$}\NormalTok{X1,}\DecValTok{30}\NormalTok{)}
\NormalTok{ech30}
\end{Highlighting}
\end{Shaded}

\begin{verbatim}
## $echantillon
##  [1] 2716 3662 2909 3068 3533 3235 3104 2462 3658 3430 3084 3750 3594 3114 3289
## [16] 2213 2837 3274 3003 2428 2871 3439 2963 2809 2771 2994 3182 2966 3080 3027
## 
## $moy_emp
## [1] 3082.167
## 
## $var_emp
## [1] -3861752
\end{verbatim}

Pour une taille d'échantillons de 60:

\begin{Shaded}
\begin{Highlighting}[]
\NormalTok{ech60}\OtherTok{=}\FunctionTok{analyse}\NormalTok{(poids}\SpecialCharTok{$}\NormalTok{X1,}\DecValTok{60}\NormalTok{)}
\NormalTok{ech60}
\end{Highlighting}
\end{Shaded}

\begin{verbatim}
## $echantillon
##  [1] 2724 3930 2815 2849 3190 3692 2890 2864 3002 3047 2945 2476 3216 3830 3294
## [16] 2943 3136 3012 3416 3642 3330 3845 3479 2780 4456 2877 3426 3468 3220 2719
## [31] 3012 4364 3046 2797 3244 4153 3028 2774 3781 3203 2771 2919 2848 2129 3181
## [46] 3059 3699 3956 3176 3257 3257 3134 3599 3487 3408 2077 3520 2961 2743 3187
## 
## $moy_emp
## [1] 3204.717
## 
## $var_emp
## [1] -12729773
\end{verbatim}

Pour une taille d'échantillons de 80:

\begin{Shaded}
\begin{Highlighting}[]
\NormalTok{ech80}\OtherTok{=}\FunctionTok{analyse}\NormalTok{(poids}\SpecialCharTok{$}\NormalTok{X1,}\DecValTok{80}\NormalTok{)}
\NormalTok{ech80}
\end{Highlighting}
\end{Shaded}

\begin{verbatim}
## $echantillon
##  [1] 3135 2233 2892 2243 3777 2774 2915 3408 2400 2832 3373 3594 4015 3376 2366
## [16] 3387 2508 2909 2709 2524 2654 3255 3302 2948 2836 2790 3375 2781 3146 2417
## [31] 3046 2941 2247 4547 2740 4000 3197 3215 2999 3622 3357 3505 2715 2784 3179
## [46] 4044 2315 3002 3173 3385 3130 3772 3527 3380 2802 2951 3083 2985 3328 3813
## [61] 3468 2380 3359 3617 3965 3236 3106 2958 2979 3275 2924 3054 3786 3507 2015
## [76] 3580 3309 3240 2901 2771
## 
## $moy_emp
## [1] 3113.6
## 
## $var_emp
## [1] -18055131
\end{verbatim}

\hypertarget{question-3}{%
\subsubsection{Question 3}\label{question-3}}

Écrire une fonction qui, à partir d'un vecteur de données normales,
calcule un intervalle de confiance de niveau de risque pour la moyenne
en supposant la variance connue. Calculer les intervalles de confiance à
95\% pour le poids moyen à la naissance à partir des échantillons
obtenus à la question.

on prend \(\sigma^2=438,9^2\)

\begin{Shaded}
\begin{Highlighting}[]
\NormalTok{interv\_conf }\OtherTok{\textless{}{-}} \ControlFlowTok{function}\NormalTok{(X,sigma,risque)\{}
\NormalTok{  n}\OtherTok{=}\FunctionTok{length}\NormalTok{(X)}
\NormalTok{  moy\_emp}\OtherTok{=}\NormalTok{ (}\DecValTok{1}\SpecialCharTok{/}\NormalTok{n)}\SpecialCharTok{*}\FunctionTok{sum}\NormalTok{(X)}
\NormalTok{  q}\OtherTok{\textless{}{-}}\FunctionTok{qnorm}\NormalTok{(}\DecValTok{1}\SpecialCharTok{{-}}\NormalTok{risque}\SpecialCharTok{/}\DecValTok{2}\NormalTok{)}
  
  \FunctionTok{return}\NormalTok{(}\FunctionTok{list}\NormalTok{(moy\_emp}\SpecialCharTok{{-}}\NormalTok{((sigma)}\SpecialCharTok{/}\FunctionTok{sqrt}\NormalTok{(n))}\SpecialCharTok{*}\NormalTok{q,moy\_emp}\SpecialCharTok{+}\NormalTok{((sigma)}\SpecialCharTok{/}\FunctionTok{sqrt}\NormalTok{(n))}\SpecialCharTok{*}\NormalTok{q))}
\NormalTok{\}}
\end{Highlighting}
\end{Shaded}

Calculon l'interval de confiance à 95\%

\begin{Shaded}
\begin{Highlighting}[]
\NormalTok{sigma}\OtherTok{=}\NormalTok{(}\FloatTok{438.9}\NormalTok{)}

\NormalTok{ic\_30}\OtherTok{=}\FunctionTok{interv\_conf}\NormalTok{(ech30}\SpecialCharTok{$}\NormalTok{echantillon,sigma,}\FloatTok{0.05}\NormalTok{)}
\NormalTok{ic\_60}\OtherTok{=}\FunctionTok{interv\_conf}\NormalTok{(ech60}\SpecialCharTok{$}\NormalTok{echantillon,sigma,}\FloatTok{0.05}\NormalTok{)}
\NormalTok{ic\_80}\OtherTok{=}\FunctionTok{interv\_conf}\NormalTok{(ech80}\SpecialCharTok{$}\NormalTok{echantillon,sigma,}\FloatTok{0.05}\NormalTok{)}


\FunctionTok{library}\NormalTok{(knitr)}

\CommentTok{\# Définir les tailles d\textquotesingle{}échantillons}
\NormalTok{tailles\_echantillons }\OtherTok{\textless{}{-}} \FunctionTok{c}\NormalTok{(}\DecValTok{30}\NormalTok{, }\DecValTok{60}\NormalTok{, }\DecValTok{80}\NormalTok{)}


\NormalTok{intervalle\_confiance }\OtherTok{\textless{}{-}} \FunctionTok{list}\NormalTok{(}\FunctionTok{c}\NormalTok{(ic\_30[}\DecValTok{1}\NormalTok{], ic\_30[}\DecValTok{2}\NormalTok{]), }\FunctionTok{c}\NormalTok{(ic\_60[}\DecValTok{1}\NormalTok{], ic\_60[}\DecValTok{2}\NormalTok{]), }\FunctionTok{c}\NormalTok{(ic\_80[}\DecValTok{1}\NormalTok{], ic\_80[}\DecValTok{2}\NormalTok{]))}

\CommentTok{\# Créer le tableau}
\NormalTok{tableau }\OtherTok{\textless{}{-}} \FunctionTok{data.frame}\NormalTok{(}\AttributeTok{Taille =}\NormalTok{ tailles\_echantillons, }\StringTok{\textasciigrave{}}\AttributeTok{Intervalle de confiance}\StringTok{\textasciigrave{}} \OtherTok{=} \FunctionTok{sapply}\NormalTok{(intervalle\_confiance, }\ControlFlowTok{function}\NormalTok{(x) }\FunctionTok{paste0}\NormalTok{(}\StringTok{"["}\NormalTok{, x[}\DecValTok{1}\NormalTok{], }\StringTok{", "}\NormalTok{, x[}\DecValTok{2}\NormalTok{], }\StringTok{"]"}\NormalTok{)))}

\CommentTok{\# Afficher le tableau avec kable}
\FunctionTok{kable}\NormalTok{(tableau)}
\end{Highlighting}
\end{Shaded}

\begin{longtable}[]{@{}rl@{}}
\toprule()
Taille & Intervalle.de.confiance \\
\midrule()
\endhead
30 & {[}2925.11120473117, 3239.22212860216{]} \\
60 & {[}3093.66168450969, 3315.77164882364{]} \\
80 & {[}3017.42356423524, 3209.77643576477{]} \\
\bottomrule()
\end{longtable}

\hypertarget{question-4}{%
\subsubsection{Question 4}\label{question-4}}

Écrire une fonction qui, à partir d'un vecteur de données normales,
calcule un intervalle de confiance de niveau de risque pour la moyenne
en supposant la variance inconnue. Calculer les intervalles de confiance
à 95\% pour le poids moyen à la naissance à partir des échantillons
obtenus à la question 2. Comparer les résultats avec ceux de la fonction
t.test.

\begin{Shaded}
\begin{Highlighting}[]
\NormalTok{interv\_conf2}\OtherTok{\textless{}{-}}\ControlFlowTok{function}\NormalTok{(X,risque)\{}
\NormalTok{  n}\OtherTok{=}\FunctionTok{length}\NormalTok{(ech30}\SpecialCharTok{$}\NormalTok{echantillon)}
\NormalTok{  moy\_emp}\OtherTok{=}\NormalTok{ (}\DecValTok{1}\SpecialCharTok{/}\NormalTok{n)}\SpecialCharTok{*}\FunctionTok{sum}\NormalTok{(ech30}\SpecialCharTok{$}\NormalTok{echantillon)}
 
\NormalTok{  var\_emp}\OtherTok{=}\NormalTok{ (}\DecValTok{1}\SpecialCharTok{/}\NormalTok{(n}\DecValTok{{-}1}\NormalTok{))}\SpecialCharTok{*}\FunctionTok{sum}\NormalTok{(((ech30}\SpecialCharTok{$}\NormalTok{echantillon)}\SpecialCharTok{{-}}\NormalTok{moy\_emp)}\SpecialCharTok{\^{}}\DecValTok{2}\NormalTok{)}
 

\NormalTok{  s}\OtherTok{\textless{}{-}}\FunctionTok{sqrt}\NormalTok{(var\_emp)}\SpecialCharTok{/}\FunctionTok{sqrt}\NormalTok{(n)}
\NormalTok{  quart }\OtherTok{\textless{}{-}} \FunctionTok{qt}\NormalTok{(}\DecValTok{1}\SpecialCharTok{{-}}\NormalTok{risque}\SpecialCharTok{/}\DecValTok{2}\NormalTok{,n}\DecValTok{{-}1}\NormalTok{)}
  \FunctionTok{return}\NormalTok{(}\FunctionTok{list}\NormalTok{(moy\_emp}\SpecialCharTok{{-}}\NormalTok{(s)}\SpecialCharTok{*}\NormalTok{quart,moy\_emp}\SpecialCharTok{+}\NormalTok{(s)}\SpecialCharTok{*}\NormalTok{quart))}
\NormalTok{\}}
\end{Highlighting}
\end{Shaded}

Calculon l'interval de confiance à 95\% et Comparons avec t.test

\begin{Shaded}
\begin{Highlighting}[]
\NormalTok{ic2\_30}\OtherTok{=}\FunctionTok{interv\_conf2}\NormalTok{(ech30}\SpecialCharTok{$}\NormalTok{echantillon,}\FloatTok{0.05}\NormalTok{)}
\NormalTok{ic3\_30}\OtherTok{=}\FunctionTok{t.test}\NormalTok{(ech30}\SpecialCharTok{$}\NormalTok{echantillon)}

\NormalTok{ic2\_60}\OtherTok{=}\FunctionTok{interv\_conf2}\NormalTok{(ech60}\SpecialCharTok{$}\NormalTok{echantillon,}\FloatTok{0.05}\NormalTok{)}
\NormalTok{ic3\_60}\OtherTok{=}\FunctionTok{t.test}\NormalTok{(ech60}\SpecialCharTok{$}\NormalTok{echantillon)}

\NormalTok{ic2\_80}\OtherTok{=}\FunctionTok{interv\_conf2}\NormalTok{(ech80}\SpecialCharTok{$}\NormalTok{echantillon,}\FloatTok{0.05}\NormalTok{)}
\NormalTok{ic3\_80}\OtherTok{=}\FunctionTok{t.test}\NormalTok{(ech80}\SpecialCharTok{$}\NormalTok{echantillon)}



\NormalTok{tableau }\OtherTok{\textless{}{-}} \FunctionTok{matrix}\NormalTok{(}\FunctionTok{c}\NormalTok{(}\FunctionTok{paste}\NormalTok{(}\StringTok{"["}\NormalTok{, ic2\_30[}\DecValTok{1}\NormalTok{], }\StringTok{","}\NormalTok{, ic2\_30[}\DecValTok{2}\NormalTok{], }\StringTok{"]"}\NormalTok{, }\AttributeTok{sep =} \StringTok{" "}\NormalTok{),}
                    \FunctionTok{paste}\NormalTok{(}\StringTok{"["}\NormalTok{, ic3\_30[}\DecValTok{4}\NormalTok{]}\SpecialCharTok{$}\NormalTok{conf.int[}\DecValTok{1}\NormalTok{],}\StringTok{","}\NormalTok{,ic3\_30[}\DecValTok{4}\NormalTok{]}\SpecialCharTok{$}\NormalTok{conf.int[}\DecValTok{2}\NormalTok{], }\StringTok{"]"}\NormalTok{, }\AttributeTok{collapse =} \StringTok{" "}\NormalTok{),}
                    \FunctionTok{paste}\NormalTok{(}\StringTok{"["}\NormalTok{, ic2\_60[}\DecValTok{1}\NormalTok{], }\StringTok{","}\NormalTok{, ic2\_60[}\DecValTok{2}\NormalTok{], }\StringTok{"]"}\NormalTok{, }\AttributeTok{sep =} \StringTok{" "}\NormalTok{),}
                    \FunctionTok{paste}\NormalTok{(}\StringTok{"["}\NormalTok{, ic3\_60[}\DecValTok{4}\NormalTok{]}\SpecialCharTok{$}\NormalTok{conf.int[}\DecValTok{1}\NormalTok{],}\StringTok{","}\NormalTok{,ic3\_60[}\DecValTok{4}\NormalTok{]}\SpecialCharTok{$}\NormalTok{conf.int[}\DecValTok{2}\NormalTok{], }\StringTok{"]"}\NormalTok{, }\AttributeTok{collapse =} \StringTok{" "}\NormalTok{),}
                    \FunctionTok{paste}\NormalTok{(}\StringTok{"["}\NormalTok{, ic2\_80[}\DecValTok{1}\NormalTok{], }\StringTok{","}\NormalTok{, ic2\_80[}\DecValTok{2}\NormalTok{], }\StringTok{"]"}\NormalTok{, }\AttributeTok{sep =} \StringTok{" "}\NormalTok{),}
                    \FunctionTok{paste}\NormalTok{(}\StringTok{"["}\NormalTok{, ic3\_80[}\DecValTok{4}\NormalTok{]}\SpecialCharTok{$}\NormalTok{conf.int[}\DecValTok{1}\NormalTok{],}\StringTok{","}\NormalTok{,ic3\_80[}\DecValTok{4}\NormalTok{]}\SpecialCharTok{$}\NormalTok{conf.int[}\DecValTok{2}\NormalTok{], }\StringTok{"]"}\NormalTok{, }\AttributeTok{collapse =} \StringTok{" "}\NormalTok{)}
\NormalTok{                    ), }
                  \AttributeTok{ncol =} \DecValTok{2}\NormalTok{, }\AttributeTok{byrow =} \ConstantTok{TRUE}\NormalTok{)}


\FunctionTok{colnames}\NormalTok{(tableau) }\OtherTok{\textless{}{-}} \FunctionTok{c}\NormalTok{(}\StringTok{"Fonction TP"}\NormalTok{, }\StringTok{"t.test"}\NormalTok{)}
\FunctionTok{rownames}\NormalTok{(tableau) }\OtherTok{\textless{}{-}} \FunctionTok{c}\NormalTok{(}\StringTok{"1"}\NormalTok{, }\StringTok{"2"}\NormalTok{, }\StringTok{"3"}\NormalTok{)}

\FunctionTok{library}\NormalTok{(knitr)}
\FunctionTok{kable}\NormalTok{(tableau, }\AttributeTok{caption =} \StringTok{"Comparaison avec t.test"}\NormalTok{)}
\end{Highlighting}
\end{Shaded}

\begin{longtable}[]{@{}
  >{\raggedright\arraybackslash}p{(\columnwidth - 2\tabcolsep) * \real{0.5000}}
  >{\raggedright\arraybackslash}p{(\columnwidth - 2\tabcolsep) * \real{0.5000}}@{}}
\caption{Comparaison avec t.test}\tabularnewline
\toprule()
\begin{minipage}[b]{\linewidth}\raggedright
Fonction TP
\end{minipage} & \begin{minipage}[b]{\linewidth}\raggedright
t.test
\end{minipage} \\
\midrule()
\endfirsthead
\toprule()
\begin{minipage}[b]{\linewidth}\raggedright
Fonction TP
\end{minipage} & \begin{minipage}[b]{\linewidth}\raggedright
t.test
\end{minipage} \\
\midrule()
\endhead
{[} 2943.57526497061 , 3220.75806836273 {]} & {[} 2943.57526497061 ,
3220.75806836273 {]} \\
{[} 2943.57526497061 , 3220.75806836273 {]} & {[} 3083.71136217004 ,
3325.7219711633 {]} \\
{[} 2943.57526497061 , 3220.75806836273 {]} & {[} 3006.54066241472 ,
3220.65933758528 {]} \\
\bottomrule()
\end{longtable}

On trouve les mêmes intervalles de confiance calculés avec notre
fonction et avec la fonction t.test

\hypertarget{question-5}{%
\subsubsection{Question 5}\label{question-5}}

Déterminer un intervalle de confiance à 95\% pour la proportion de
cartons de lait défectueux en utilisant les formules vues en cours et en
utilisant la fonction binom.test de R.

On calcul la la proportion de carton de lait défectueux avec les
formules vues en cours

Les formules sont:
\[\Bigg[\ \bar{X_n} - u_{1-\frac\alpha 2} \sqrt{ \frac{\bar{X_n}(1-\bar{X_n})}{n}}\ ,\ \bar{X_n} + u_{1-\frac\alpha 2} \sqrt{ \frac{\bar{X_n}(1-\bar{X_n})}{n}}\ \Bigg]\]

\begin{Shaded}
\begin{Highlighting}[]
\NormalTok{prop}\OtherTok{=}\DecValTok{20}\SpecialCharTok{/}\DecValTok{197}
\NormalTok{r}\OtherTok{=}\FloatTok{0.05}

\NormalTok{prop\_inf}\OtherTok{=}\NormalTok{ prop}\SpecialCharTok{{-}}\NormalTok{ (}\FunctionTok{qnorm}\NormalTok{(}\DecValTok{1}\SpecialCharTok{{-}}\NormalTok{(r}\SpecialCharTok{/}\DecValTok{2}\NormalTok{),}\AttributeTok{mean=}\DecValTok{0}\NormalTok{,}\AttributeTok{sd=}\DecValTok{1}\NormalTok{)}\SpecialCharTok{*}\FunctionTok{sqrt}\NormalTok{(prop}\SpecialCharTok{*}\NormalTok{(}\DecValTok{1}\SpecialCharTok{{-}}\NormalTok{prop)}\SpecialCharTok{/}\DecValTok{197}\NormalTok{))}

\NormalTok{prop\_sup}\OtherTok{=}\NormalTok{prop }\SpecialCharTok{+}\NormalTok{ (}\FunctionTok{qnorm}\NormalTok{(}\DecValTok{1}\SpecialCharTok{{-}}\NormalTok{(r}\SpecialCharTok{/}\DecValTok{2}\NormalTok{),}\AttributeTok{mean=}\DecValTok{0}\NormalTok{,}\AttributeTok{sd=}\DecValTok{1}\NormalTok{)}\SpecialCharTok{*}\FunctionTok{sqrt}\NormalTok{(prop}\SpecialCharTok{*}\NormalTok{(}\DecValTok{1}\SpecialCharTok{{-}}\NormalTok{prop)}\SpecialCharTok{/}\DecValTok{197}\NormalTok{))}
\end{Highlighting}
\end{Shaded}

On calcul la la proportion de carton de lait défectueux avec la fonction
binom.test

\begin{Shaded}
\begin{Highlighting}[]
\NormalTok{interv\_conf}\OtherTok{\textless{}{-}}\FunctionTok{binom.test}\NormalTok{(}\FunctionTok{c}\NormalTok{(}\DecValTok{20}\NormalTok{,}\DecValTok{197}\NormalTok{),}\AttributeTok{n=}\DecValTok{197}\NormalTok{,}\AttributeTok{conf.level =} \DecValTok{1}\SpecialCharTok{{-}}\NormalTok{r)}
\end{Highlighting}
\end{Shaded}

\begin{Shaded}
\begin{Highlighting}[]
\NormalTok{dif1}\OtherTok{=}\FunctionTok{abs}\NormalTok{((prop\_inf}\SpecialCharTok{{-}}\NormalTok{interv\_conf[}\DecValTok{4}\NormalTok{]}\SpecialCharTok{$}\NormalTok{conf.int[}\DecValTok{1}\NormalTok{])}\SpecialCharTok{/}\NormalTok{interv\_conf[}\DecValTok{4}\NormalTok{]}\SpecialCharTok{$}\NormalTok{conf.int[}\DecValTok{1}\NormalTok{])}\SpecialCharTok{*}\DecValTok{100}
\NormalTok{dif2}\OtherTok{=}\FunctionTok{abs}\NormalTok{((prop\_sup}\SpecialCharTok{{-}}\NormalTok{interv\_conf[}\DecValTok{4}\NormalTok{]}\SpecialCharTok{$}\NormalTok{conf.int[}\DecValTok{2}\NormalTok{])}\SpecialCharTok{/}\NormalTok{interv\_conf[}\DecValTok{4}\NormalTok{]}\SpecialCharTok{$}\NormalTok{conf.int[}\DecValTok{2}\NormalTok{])}\SpecialCharTok{*}\DecValTok{100}

\FunctionTok{library}\NormalTok{(knitr)}


\CommentTok{\# Créer le tableau}
\NormalTok{tableau }\OtherTok{\textless{}{-}} \FunctionTok{data.frame}\NormalTok{(}
  \StringTok{\textasciigrave{}}\AttributeTok{Interval formules du cours}\StringTok{\textasciigrave{}} \OtherTok{=} \FunctionTok{paste}\NormalTok{(}\StringTok{"["}\NormalTok{, prop\_inf, }\StringTok{","}\NormalTok{, prop\_sup, }\StringTok{"]"}\NormalTok{, }\AttributeTok{sep =} \StringTok{""}\NormalTok{),}
  \StringTok{\textasciigrave{}}\AttributeTok{Interval binom.test}\StringTok{\textasciigrave{}} \OtherTok{=} \FunctionTok{paste}\NormalTok{(}\StringTok{"["}\NormalTok{, interv\_conf[}\DecValTok{4}\NormalTok{]}\SpecialCharTok{$}\NormalTok{conf.int[}\DecValTok{1}\NormalTok{], }\StringTok{","}\NormalTok{, interv\_conf[}\DecValTok{4}\NormalTok{]}\SpecialCharTok{$}\NormalTok{conf.int[}\DecValTok{2}\NormalTok{], }\StringTok{"]"}\NormalTok{, }\AttributeTok{sep =} \StringTok{""}\NormalTok{),}
  \StringTok{\textasciigrave{}}\AttributeTok{Ecart inférieur}\StringTok{\textasciigrave{}} \OtherTok{=}\NormalTok{ dif1,}
  \StringTok{\textasciigrave{}}\AttributeTok{Ecart supérieur}\StringTok{\textasciigrave{}} \OtherTok{=}\NormalTok{ dif2}
\NormalTok{)}

\CommentTok{\# Afficher le tableau avec kable}
\FunctionTok{kable}\NormalTok{(tableau)}
\end{Highlighting}
\end{Shaded}

\begin{longtable}[]{@{}
  >{\raggedright\arraybackslash}p{(\columnwidth - 6\tabcolsep) * \real{0.3578}}
  >{\raggedright\arraybackslash}p{(\columnwidth - 6\tabcolsep) * \real{0.3486}}
  >{\raggedleft\arraybackslash}p{(\columnwidth - 6\tabcolsep) * \real{0.1468}}
  >{\raggedleft\arraybackslash}p{(\columnwidth - 6\tabcolsep) * \real{0.1468}}@{}}
\toprule()
\begin{minipage}[b]{\linewidth}\raggedright
Interval.formules.du.cours
\end{minipage} & \begin{minipage}[b]{\linewidth}\raggedright
Interval.binom.test
\end{minipage} & \begin{minipage}[b]{\linewidth}\raggedleft
Ecart.inférieur
\end{minipage} & \begin{minipage}[b]{\linewidth}\raggedleft
Ecart.supérieur
\end{minipage} \\
\midrule()
\endhead
{[}0.0593482998838742,0.143697385395314{]} &
{[}0.0572091306283455,0.13875216342604{]} & 3.739209 & 3.564068 \\
\bottomrule()
\end{longtable}

On remarque des écart de 3.5\% et 3.7\% ce qui est acceptable mais quand
même présent.

\end{document}
